\documentclass{article}
\begin{document}
\section*{Vertical arrows and the evaluation constraints}
The strict fibre of the Grothendieck projection is not replaced by its set of objects. A vertical morphism carries a proof of lifting the base identity; its decoded payload arrow must be corrected by the equality induced by that proof. The canonical inclusion and explicit decoder compose to the identity as functors in both orders. The square comparing a chosen pull with the actual indexed morphism functor is invertible and natural on all vertical arrows.
\par For a positive native power and fixed ordinary payload this gives a pointwise equivalence to the actual iterate value and a square for each base map. In the strict-fibre presentation, two further equations compare the actual evaluation squares for an identity and a composite. Both constraints occur with their proper orientation: for a composite, following its square by the target compositor equals the source compositor followed by the two horizontally pasted squares. The proof cancels against the final cartesian lift; it does not replace a compositor by an identity.
\par Not proved here: independent constraint equations for the pasted value-target squares, mixed base/payload coherence or a bundled presheaf equivalence. The positive length convention is unchanged and includes neither the empty summand nor flattening or strict algebra recovery.
\end{document}
